%% 自動生成（tools/mknext.py）。各章の道しるべ「この章の先に」から、後ろの巻の章を集めたもの
\begin{longtblr}{colspec={Q[l,m,wd=58mm] X[l,m]}, width=\linewidth, rowhead=1, row{1}={bg=bkVol!14, font=\sffamily\bfseries\footnotesize},
  hlines={0.4pt,bkLine}, rowsep=1.5pt, cells={font=\footnotesize}}
この巻の章 & この道具を使う章 \\
{\bknextpill{4-1}\ きまりから考えを積み上げる：定義・公理・定理} & {\bknextpill{6-15}\ まとめ：物理は「微分方程式＋はじめの状態」で書ける} \\
{\bknextpill{4-2}\ 角度と円、ラジアン} & {\bknextpill{5-3}\ かぎりなく近づく：極限と接線\\ \bknextpill{5-15}\ sin・cosを微分すると？\\ \bknextpill{5-17}\ 近くでは多項式で近似できる：テイラー展開} \\
{\bknextpill{4-3}\ にた形と比：相似} & {\bknextpill{5-10}\ 細い長方形に分けてたそう} \\
{\bknextpill{4-4}\ 三平方の定理} & {\bknextpill{5-16}\ 積分の道具箱：置換積分と部分積分} \\
{\bknextpill{4-5}\ 面積と体積：細かく分けてたす} & {\bknextpill{5-4}\ 表から導関数を見つけよう\\ \bknextpill{5-6}\ 微分の道具箱：積と合成関数\\ \bknextpill{5-9}\ グラフの面積は進んだ距離\\ \bknextpill{5-10}\ 細い長方形に分けてたそう} \\
{\bknextpill{4-6}\ 坂道と三角比} & {\bknextpill{5-17}\ 近くでは多項式で近似できる：テイラー展開} \\
{\bknextpill{4-8}\ 三角関数：円運動の影} & {\bknextpill{5-15}\ sin・cosを微分すると？\\ \bknextpill{5-16}\ 積分の道具箱：置換積分と部分積分\\ \bknextpill{6-6}\ ばねと振り子\\ \bknextpill{7-6}\ 回転行列：くるっと回す} \\
{\bknextpill{4-9}\ 波をたし合わせる：加法定理とうなり} & {\bknextpill{5-15}\ sin・cosを微分すると？\\ \bknextpill{6-10}\ オイラーの公式\\ \bknextpill{9-1}\ 波をたし合わせて形を作る：フーリエ級数} \\
{\bknextpill{4-10}\ 向きと大きさをもつ量：ベクトル} & {\bknextpill{7-2}\ 行列×ベクトル：形を変える機械\\ \bknextpill{8-5}\ ベクトル場} \\
{\bknextpill{4-11}\ ベクトルを成分で計算しよう} & {\bknextpill{5-15}\ sin・cosを微分すると？\\ \bknextpill{7-2}\ 行列×ベクトル：形を変える機械} \\
{\bknextpill{4-12}\ 内積：どれだけ同じ向きか} & {\bknextpill{5-12}\ 物理での積分：仕事と流量\\ \bknextpill{8-7}\ 道にそってたす：線積分\\ \bknextpill{9-2}\ 成分の取り出し方：直交性} \\
{\bknextpill{4-13}\ 空間の座標と3次元のベクトル} & {\bknextpill{7-7}\ 3次元の回転と座標の変換\\ \bknextpill{8-5}\ ベクトル場} \\
{\bknextpill{4-14}\ 外積：回す力と右ねじ} & {\bknextpill{7-4}\ 行列式：面積が何倍になるか\\ \bknextpill{8-10}\ うず：回転 rot} \\
{\bknextpill{4-15}\ いろいろな座標：極座標・円柱座標・球座標} & {\bknextpill{8-6}\ 面や立体でたし合わせる：重積分\\ \bknextpill{8-12}\ ならす働き：ラプラシアン} \\
\end{longtblr}
